\subsection{Parametry przepływu i własności fizyczne powietrza}
stanu gazu doskonale idealnego:
\begin{equation}
    \rho = \frac{p}{R \cdot T}
\end{equation}
gdzie:
\begin{itemize}
    \item $p$ - ciśnienie atmosferyczne,
    \item $R$ - stała gazowa dla powietrza $R \approx 287 \, \text{J}/(\text{kg}\cdot\text{K})$,
    \item $T$ - temperatura bezwzględna w kelwinach.
\end{itemize}

Znając gęstość powietrza oraz ciśnienie dynamiczne $q$ zmierzone w tunelu, prędkość przepływu niezakłóconego $V$ oblicza się, przekształcając wzór na ciśnienie dynamiczne:
\begin{equation}
    q = \frac{\rho V^2}{2} \implies V = \sqrt{\frac{2q}{\rho}}
\end{equation}

Z kolei kluczowym bezwymiarowym kryterium podobieństwa charakteryzującym opływ jest liczba Reynoldsa $Re$, definiowana za pomocą cięciwy profilu $c$, prędkości $V$ oraz lepkości kinematycznej powietrza $\nu$:
\begin{equation}
    Re = \frac{V \cdot c}{\nu}
\end{equation}

\subsection{Współczynnik siły nośnej}
W celu opisu wektora siły aerodynamicznej w układzie współrzędnych, rzut siły na oś prostopadłą do kierunku prędkości definiuje siłę nośną $P_z$. Aby uniezależnić wyniki od skali modelu i warunków pomiaru, siłę nośną opisuje się za pomocą bezwymiarowego współczynnika $C_z$:
\begin{equation}
    C_z = \frac{P_z}{\frac{\rho V^2}{2}S} = \frac{P_z}{qS}
\end{equation}
gdzie $S$ oznacza powierzchnię odniesienia. Dla ustalonej wartości liczby Reynoldsa współczynnik ten zależy wyłącznie od kąta natarcia $\alpha$.
